\section{Gaza}
\label{sec:6.8}
In Gaza the term applies on both sides, and the licenses attach to acts. The case also rests on both versions of the term, so each step below says which.

\emph{7 October: Israel's license} (the Charter as courts read it). On 7 October Israel's license attached for force repelling the attack: engaging the attackers inside Israel, and the launchers and staging areas conducting it, for as long as fire continued and was corroborated, which it was, openly, for months.\note{
That license rests not on Article 51, which the Wall opinion held irrelevant to attacks originating in territory a state controls (para.~139), a holding whose reach to Gaza after 2005 is disputed, but on the duty the force term names, which the same opinion affirmed (para.~141) and which no dispute about control touches.
}
It licenses stopping the massacre and nothing beyond it.

\emph{What that license never covered} (person-directed terms; for ranking and timing, the machine-selection departure). That license covered the least force that repels. It never covered a campaign to destroy Hamas, a ground occupation, a siege of the population, the ranking of people by affiliation, or the timing of strikes to family presence.\note{
Those fall outside it on proportionality and are forbidden by the person-directed terms, from the first week. Ranking by affiliation and timing to family presence are barred by the machine-selection departure; the siege and the disproportionate campaign are barred without one. A machine in the IDF pipeline would refuse them from October 2023.
}

\emph{The attack itself} (person-directed terms, and the rule that an occupation is not an attack under way). A machine ordered by Hamas to assist the attack of 7 October refuses too.\note{
An occupation is not an attack under way, so nothing that morning was repelling anything, and the acts themselves, killing civilians and taking hostages, are forbidden by the person-directed terms whatever the ad bellum position.
}

\emph{After 7 October: the Palestinian license} (extended term: standing to defend). After 7 October the Gaza population is a defender's population: the territory is occupied on a disinterested finding, and Palestine holds observer-state status by a vote of 138 to 9. A machine deployed by Palestinian forces holds an emergency license to repel attacks under way on that population, such as engaging a unit conducting a massacre or intercepting a munition aimed at a shelter.\note{
It holds none to defend launcher crews, who are the attacking force Israel's license covers, and none for fire on Israeli cities.
}

\emph{September 2025: the genocide finding} (extended term: the genocide exception). The September 2025 finding \autocite{coiopt2025genocide}, taken with the Court's provisional-measures order of January 2024 \autocite{icj2024sagenocidepm}, meets the two-body rule.\note{
Israel and the United States dispute the finding \autocite{cnn2025coi}, and the Court's merits judgment is years off, with written pleadings scheduled to close in 2029 \autocite{icj2026sagenocide}; the two-body rule is what stands in for it until then.
}
It withdraws whatever license Israeli forces still held for the force and place it covers, and opens the exception for force against those forces there.

The finding changes which exception each side's force rests on. It doesn't change what the person-directed terms forbade all along.

Run the term's own clock on the case and the cost of corroboration shows in dates. Attacks on the Palestinian population were covered by the emergency license from 1948, on corroboration that existed: the mediator's report \autocite{bernadotte1948progress} and the Red Cross delegate's account \autocite{dereynier1950jerusalem}. Standing to defend the territory attached by the end of 1967, on the Council's demand for withdrawal \autocite{unsc1967res242,unsc1979res446}. Force to end the occupation attached for the West Bank in 2004, on the Court's Wall opinion \autocite{icj2004wall}, and for Gaza in 2009, on the fact-finding mission \autocite{hrc2009gaza}, each paired with that demand, and the Court's 2024 opinion closed the question by the single-court route \autocite{icj2024palestine}. The genocide exception needed the Court's order of January 2024 and the commission's finding of September 2025. So the term licensed defense of the people throughout, defense of the territory from the occupation's first year, and force against the genocide after nearly two years of it, on a schedule the occupied didn't set; and on 7 October it licensed Israel to stop the massacre for as long as the fire continued and was corroborated, and refused the attack itself and every launcher crew on the other side. Every one of those dates is a document, and a reader who thinks the term was drawn to reach this result has one move: to say which of the documents shouldn't count, a Council resolution, a fact-finding mission, a commission of inquiry, an advisory opinion. Appendix~\ref{sec:B.1} says what each is worth and why, and the term reads them the same way for every war in this book. What the dates leave standing is the wait, thirty-seven years to the first finding by a body that took evidence and twenty months from the Court's order to the commission, and the reply is the one the term is built on: a faster rule lifts on the intervener's own claim, and that is Gleiwitz.
